\chapter{总结与展望}

\section{论文工作总结}

% TODO: 系统总结全文工作与三大创新点
%
% 本文围绕复杂城市环境下非法低轨卫星终端的高精度定位与高并发识别需求，
% 以低成本无人机为感知平台，深入研究了电磁信道建模、深度学习特征提取、
% 智能决策融合及跨模态目标检测等核心问题。主要工作和贡献总结如下：
%
% （1）提出了基于射线追踪的高保真城市信道建模与深度RF指纹定位方法...
%     关键指标：定位精度8.3m，F1-Score=0.85，8×8 MIMO在-15dB下仍<20m
%
% （2）提出了PA-DQN物理感知深度强化学习融合框架...
%     关键指标：RMSE=2.15m，vs EKF提升70.7%，vs Q-Learning提升63.6%
%
% （3）提出了基于信号成像与Anchor-Free YOLOv7的多目标并行定位方法...
%     关键指标：O(1)复杂度，~25ms/帧，置信度>0.98

\section{未来研究展望}

% TODO: 未来改进方向
%
% （1）引入DDPG等Actor-Critic架构，实现融合权重的连续调节
% （2）结合LSTM利用信道特征的时序信息，提升动态环境预判能力
% （3）扩展更多差异化城市场景验证模型泛化能力
% （4）各仿真模块的系统级联调
% （5）向实物测试与实际部署推进
